\section{Certifying Ground-Truth Depth}\label{sec:gt}
A model enters the analysis only if it passes all of the following tests:
\begin{itemize}
\item clean bit accuracy $\ge 0.99$;
\item confound-free interchange accuracy $\ge 0.98$ at every hop, on flipped bits, with denominators reported;
\item strictness $\ge 0.98$ for every non-circuit component;
\item bypass-freeness $\ge 0.98$, both trained and held-out (mean ablation);
\item circuit-only accuracy $\ge 0.99$ with all non-circuit components jointly resampled;
\item every hop-hosting MLP necessary in the circuit-only model (accuracy $\le 0.9$ when mean-ablated);
\item site separation $\ge 3$.
\end{itemize}
Table~\ref{tab:cert} summarises the results (\NCertified{} of \NModels{} models certified; by depth $D=1..5$: \NCertOne, \NCertTwo, \NCertThree, \NCertFour, \NCertFive{} of 3). Models deeper than $D=3$ need longer training. Our pre-specified rule retrains uncertified models with twice the steps.
\paragraph{Two lessons.} (1) For the Round-1 2-hop task ($y=(a\oplus b)\wedge c$), zeroing the last MLP left accuracy at $100\%$. The AND is a linear threshold that the unembedding computes, so that model has nonlinear depth 1, and a shallow graph would have been \emph{correct}. (2) Single-component strictness is not enough. The non-circuit outputs stayed large and input-dependent, and resampling several of them at once broke the circuit. This is why we add joint resampling and site reservation.
\begin{table}[t]\centering\caption{Ground-truth certificates (pilot).}\label{tab:cert}\begin{tabular}{rrrrrrr}\toprule
$D$ & models & certified & clean acc & min IIA (cf) & min strict & min bypass-free\\\midrule
1 & 5 & 5 & 1.000 & 1.000 & 1.000 & 1.000\\
2 & 5 & 5 & 1.000 & 1.000 & 1.000 & 1.000\\
3 & 5 & 4 & 1.000 & 1.000 & 0.494 & 1.000\\
4 & 5 & 1 & 1.000 & 0.994 & 0.611 & 0.999\\
5 & 5 & 0 & 1.000 & 1.000 & 0.560 & 0.997\\
\bottomrule\end{tabular}
\end{table}
